% -----------------------------------------------------------------------------------------------------
\section{状态表与开放问题}\label{sec:front-status}

\cref{tab:front-status}逐条列出本章带标签陈述的等级、最弱的一步与误差类型（\cref{def:intro-error-types}；关于单个积分、模型或显式计算的陈述记为“—”）。两条编号假设见\cref{asm:front-mor-complete,asm:front-scet-modes}，两条数学引理（\cref{lem:front-min-norm,lem:front-jantzen}）不列入表中。

\begin{table}[htbp]
\centering
\caption{第~24~章状态表（上）。“审稿”一列标出证明或结果仅见于预印本者。}\label{tab:front-status}
\footnotesize
\renewcommand{\arraystretch}{1.08}
\begin{tabularx}{\textwidth}{@{}p{1.75cm}>{\raggedright\arraybackslash}Xp{1.35cm}>{\raggedright\arraybackslash}p{3.6cm}p{0.75cm}p{0.95cm}@{}}
\toprule
标签 & 陈述 & 等级 & 最弱的一步 & 误差 & 审稿 \\
\midrule
\cref{thm:front-collins-pinch} & 物理区钳夹 $\Leftrightarrow$ Landau 点 & [A] & 证明只见于预印本；符号条件 & — & 预印本 \\
\cref{cor:front-coleman-norton} & 动量空间 Coleman--Norton 图像 & [R, $R_{\mathrm{cond}}$] & 前提\cref{thm:front-collins-pinch}（内容 [A]） & — & 预印本 \\
\cref{prop:front-anomalous} & 正一阶位移不是必要条件 & [V] & 数学反例 & — & 预印本 \\
\cref{prop:front-landau-not-singular} & Landau 点 $\not\Rightarrow$ 奇点 & [V] & 显式反例 & — & 预印本 \\
\cref{prop:front-ir-forest} & 宽角振幅红外森林公式 & [A] & 多胶子 Ward 恒等式（\cref{lem:method-ward-multi}） & E-S & 是 \\
\cref{prop:front-coordinate-space} & 坐标空间／算符路线的振幅因子化 & [A] & 只含末态共线方向；(a) 据摘要 & E-S & 是 \\
\cref{prop:front-local-ir} & 两圈局域红外因子化 & [P] & 只限所列过程与两圈 & — & 部分 \\
\cref{prop:front-ope-scalar} & $\varphi^4_4$ 中 OPE 收敛、结合律 & [P] & 标量模型 & — & 是 \\
\cref{prop:front-ope-ym} & 纯 YM 中 OPE 在规范不变算符上闭合 & [P] & 不含夸克；(b) 的界只对纯 YM & — & 部分 \\
\cref{prop:front-flow-renorm} & 流方程可重整性证明 & [P] & 无含夸克 QCD 的证明 & — & 部分 \\
\cref{conj:front-ope-convergence} & YM／QCD 中 OPE 收敛 & [C] & 需要更精细的界 & — & — \\
\cref{prop:front-lp-facets} & 系数全正时的面展开 & [R, $R_{\mathrm{PT}}$] & 多面、全阶情形为猜想 & — & 是 \\
\cref{prop:front-onshell-facets} & 在壳展开的面判据 & [R, $R_{\mathrm{PT}}$] & 完备性见\cref{asm:front-mor-complete} & — & 是 \\
\cref{prop:front-commuting} & 成对软区域与展开可交换 & [R, $R_{\mathrm{cond}}$] & \cref{asm:front-mor-complete}（[C]） & — & 是 \\
\cref{prop:front-hidden-regions} & “区域 $=$ 下表面”在闵氏运动学 & [V] & 三圈非平面反例 & — & 是 \\
\cref{conj:front-mor-minkowski} & 闵氏极限下区域方法有效 & [C] & 无一般论证 & — & — \\
\cref{prop:front-glauber-bfkl} & Glauber 有效理论、BFKL & [A] & (i) 是论证；“far from a proof” & E-S & 是 \\
\cref{prop:front-syz} & 纯 Glauber 梯形图相消 & [A] & 只限纯梯形图、单标度观测量 & E-S & 是 \\
\cref{conj:front-mpi-violation} & MPI 敏感观测量的因子化破坏 & [C] & Gaunt 论证的纯 Glauber 部分已被否定 & — & — \\
\cref{prop:front-heavy-fsi} & 重末态的软末态相互作用 & [A] & 据摘要 & E-S & 是 \\
\cref{prop:front-sll-all-orders} & $\mathcal H\otimes\mathcal W$ 与 SLL 的 LL 重求和 & [A] & 有效理论推导，单圈重整化群 & E-S & 是 \\
\cref{prop:front-sll-closed-form} & SLL 塔与单态道闭式 & [A] & 继承上一条 & E-S & 是 \\
\cref{prop:front-restoration} & 间隙截面的因子化恢复 & [P] & 三圈、LL & — & 是 \\
\cref{prop:front-jettiness-cvl} & 1-jettiness 标准指数化 & [V] & LL 相干性破坏计算 & — & 是 \\
\cref{thm:front-hgg} & $h\to\gamma\gamma$（$b$ 圈）端点因子化 & [A] & \cref{asm:front-scet-modes} & E-S & 是 \\
\cref{prop:front-ggh} & 轻夸克 $gg\to h$、RBS 方案 & [A] & \cref{asm:front-scet-modes}；据摘要 & E-S & 是 \\
\cref{prop:front-nlp-threshold} & NLP 阈值 DY 的 $d$ 维因子化 & [P] & 只在 $d$ 维（裸）成立 & — & 是 \\
\cref{prop:front-nlp-ll} & NLP 的 LL 重求和与端点条件 & [A] & 只到 LL；\cref{asm:front-scet-modes} & E-S & 是 \\
\cref{prop:front-nlp-dis} & NLP DIS 与阈值 PDF & [A] & 只到 LL；\cref{asm:front-scet-modes} & E-S & 预印本 \\
\cref{conj:front-nlp-beyond-ll} & 超出 LL 的端点因子化 & [C] & 无一般定理 & — & — \\
\bottomrule
\end{tabularx}
\end{table}

\begin{table}[htbp]
\centering
\caption{第~24~章状态表（续）。}\label{tab:front-status-b}
\footnotesize
\renewcommand{\arraystretch}{1.08}
\begin{tabularx}{\textwidth}{@{}p{1.75cm}>{\raggedright\arraybackslash}Xp{1.35cm}>{\raggedright\arraybackslash}p{3.6cm}p{0.75cm}p{0.95cm}@{}}
\toprule
标签 & 陈述 & 等级 & 最弱的一步 & 误差 & 审稿 \\
\midrule
\cref{conj:front-glauber-nlp} & NLP TMD 中的 Glauber 相消 & [C] & 处处被假定 & — & — \\
\cref{prop:front-tmd-nlp-operator} & TMD 算符展开与运动学修正 & [A] & NLO；强子层次需\cref{conj:front-glauber-nlp} & E-S & 是 \\
\cref{prop:front-sidis-nlp} & SIDIS 的 NLP SCET 因子化 & [A] & \cref{asm:front-scet-modes}(SM1)；Glauber 另需猜想 & E-S & 是 \\
\cref{prop:front-nlp-tmd-oneloop} & 单圈 NLP 系数与扭度 3 演化 & [P] & 单圈；各组尚不一致 & — & 部分 \\
\cref{prop:front-msbar-negative} & $\MSbar$ PDF 可为负（(a) [R, $R_{\mathrm{PT}}$]） & [V] & 模型反例 & — & 是 \\
\cref{prop:front-positivity-large-scale} & 足够大标度上的正定性 & [P] & NLO 论证，(PS1)--(PS3) & — & 是 \\
\cref{prop:front-brst-bilocal} & 双局域 PDF 算符的矩约化 & [A] & \cref{prop:appL-renormalized}证明的循环性未核查 & — & 是 \\
\cref{conj:front-brst-multilocal} & 多局域光线算符的 BRST 分类 & [C] & 无证明（Collins 假定） & — & — \\
\cref{prop:front-nlo-bk} & NLO BK/JIMWLK、三圈 BK 核 & [A] & 程函框架内、固定对数精度 & — & 是 \\
\cref{prop:front-bk-collinear} & NLO BK 的共线重求和 & [P] & 处方不唯一 & — & 是 \\
\cref{prop:front-regge} & Wilson 线方法与 Regge 极限 & [A] & 程函框架内；(a) 只在平面 $\mathcal N=4$ 精确核对 & — & 是 \\
\cref{prop:front-regge-pole} & 简单 Regge 极点因子化（NNLL） & [V] & 基于红外因子化的计算 & — & 是 \\
\cref{prop:front-smallx-tmd} & 小 $x$ TMD 与 Sudakov 对数 & [A] & LO／LL；(c) 为 [P] & — & 是 \\
\cref{prop:front-smallx-helicity} & 小 $x$ 螺旋度截距 & [A] & 双对数、大 $N_c$；与 BER 的分歧未决 & — & 是 \\
\cref{prop:front-dps-scheme} & 无重复计数的 DPS & [A] & 固定阶的方案自洽性 & E-S & 是 \\
\cref{prop:front-rapidity-renorm} & 多部分子软因子快度重整化 & [A] & 据摘要 & — & 是 \\
\cref{prop:front-dpd-evolution} & 色关联 DPD 的演化 & [P] & 算到 NLO & — & 是 \\
\cref{prop:front-ngl} & 超出 LL 的非全局对数 & [A] & (b) 限大 $N_c$、NLL & E-S & 部分 \\
\cref{prop:front-jet-substructure} & 喷注子结构与半单举喷注函数 & [A] & (c) 为 NLO 加相容性 & E-S & 是 \\
\cref{prop:front-lightray-ope} & 能量关联子、光线 OPE & [A] & (b) 是共形场论陈述 & E-S & 是 \\
\cref{prop:front-nrqcd-ir} & NNLO 的规范完备八重态算符 & [P] & 两圈 & — & 是 \\
\cref{thm:front-luxqed} & 由结构函数得到光子 PDF & [A] & 辅助过程的共线因子化 & E-S & 是 \\
\cref{prop:front-lepton-pdf} & 轻子／电子 PDF、混合核 & [A] & 据摘要；(c) 有限阶 & E-S & 是 \\
\cref{prop:front-ew-pdf} & 电弱 PDF 与电弱对数 & [A] & 据摘要 & E-S & 是 \\
\cref{prop:front-ew-violation} & 电弱 BN 抵消失效 & [V] & 显式微扰计算 & — & 是 \\
\bottomrule
\end{tabularx}
\end{table}

\begin{tcolorbox}[enhanced,breakable,title=本章澄清的常见误解,fonttitle=\bfseries,
  colback=gray!3!white,colframe=gray!55!black,boxrule=0.6pt,arc=2pt,
  left=5pt,right=5pt,top=3pt,bottom=3pt,before skip=6pt,after skip=6pt]
\textbf{“区域方法可以确认领头区域”}（C5）：在 Minkowski 运动学中这是循环论证，区域的识别必须来自钳夹分析与幂次计数（\cref{prop:front-hidden-regions}）。\textbf{“区域方法是定理”}（M19、O11）：欧氏极限 [R]，Minkowski 极限没有一般证明，可能需要解析正规子（\cref{conj:front-mor-minkowski}）；“区域 $=$ Newton 多面体的下表面”在 Minkowski 运动学中不成立。\textbf{“Landau 方程给出奇点”}（M20、O12）：它只是奇点的必要条件；外部不变量的奇点与固定外动量下的红外发散要区分，端点与第二类奇点另需处理（\cref{prop:front-landau-not-singular}）。\textbf{“Collins（2020）的 Landau 定理已正式发表”}：截至 2026 年 10 月只有预印本（\cref{thm:front-collins-pinch}）。\textbf{“PDF 是概率密度，所以 $\MSbar$ PDF 为正”}（M2）：见\cref{prop:front-msbar-negative}。\textbf{“严格共线因子化的破坏说明 PDF 因子化失效”}（M21）：那是振幅层次的陈述（\cref{rem:front-glauber-status}），间隙截面中甚至有 PDF 因子化恢复的证据（\cref{prop:front-restoration}）。
\end{tcolorbox}

\begin{remark}[开放问题]\label{rem:front-open}
(1)~钳夹理论：去掉\cref{thm:front-collins-pinch}的符号条件与二次限制；经审稿的版本；在钳夹处确定哪些为零的分母真正钳住围道及其余分母的容许形变锥（第九章每一个近似都需要它）；系统的坐标空间版本。(2)~Minkowski 运动学中的区域方法：一个证明（\cref{conj:front-mor-minkowski}；Jantzen 的收敛步骤），以及保证找到隐藏区域（Glauber、多硬子图）的算法。(3)~强子—强子碰撞中的 Glauber 与因子化：单举 Drell--Yan 的 [R] 级证明；纯梯形图之外 MPI 敏感观测量的命运（\cref{conj:front-mpi-violation}）；三圈以上非全局观测量的 PDF 因子化；$N$-jettiness 中 $\alpha_s^nL^{2n-2}$ 相干性破坏对数的确认与重求和（及其对 $N$-jettiness 切片的影响）。(4)~次领头幂次：有入射强子时超出领头对数的端点因子化与重因子化条件的一般定理（\cref{conj:front-nlp-beyond-ll}）；NLP 的 Glauber 相消（\cref{conj:front-glauber-nlp}）；各组 NLP TMD 结果的调和。(5)~BRST：多局域分类（\cref{conj:front-brst-multilocal}）。(6)~OPE：YM／QCD 中的收敛（\cref{conj:front-ope-convergence}）；含夸克 QCD 的自洽流方程证明。(7)~正定性：从第一原理确定 $\MSbar$ PDF 转为正的标度。(8)~小 $x$：Kovchegov 等与 BER 分歧（从四圈起）的来源；NLO BK 与方案无关的共线改进；小 $x$ 与共线因子化的自洽匹配。(9)~双部分子散射：带色末态的因子化（\cref{conj:break-dps-coloured}）；色关联 DPD 的唯象。(10)~重夸克偶素：超出 NNLO 的 NRQCD 产生因子化；两圈确认背后的“一般原理”。(11)~电弱：$Q\gg M_W$ 处不因子化的电弱对数的重求和。
\end{remark}
